% DockPrepBench supplementary information. Build: latexmk supplementary.tex (see README.md).
% Every table body under tables/ is generated by `dockprepbench report`; edit captions here,
% never the numbers.
\documentclass[11pt]{article}

% Packages, units and macros shared by main.tex and supplementary.tex.
\usepackage[T1]{fontenc}
\usepackage[utf8]{inputenc}
\usepackage{lmodern}
\usepackage[margin=2.5cm]{geometry}
\usepackage{setspace}
\usepackage{amsmath}
\usepackage{graphicx}
\usepackage{booktabs}
\usepackage{longtable}
\usepackage{siunitx}
\usepackage[dvipsnames]{xcolor}
\usepackage[numbers,sort&compress]{natbib}
\usepackage{lineno}
\usepackage[hidelinks]{hyperref}

\graphicspath{{../docs/figures/}{../figures/}}
\sisetup{detect-all}
% Ångström is no longer an SI unit, so siunitx deprecates it; it remains standard here.
\DeclareSIUnit{\angstrom}{\text{\AA}}
\onehalfspacing
\linenumbers

% Placeholders that stay visible until filled. Search for \result and \todo before submitting.
\newcommand{\result}[1]{\textcolor{orange}{[#1]}}
\newcommand{\todo}[1]{\textcolor{red}{[TODO: #1]}}

% Notation used throughout. (siunitx already defines \ang, for angles.)
\newcommand{\pbvalid}{PB-valid}


\renewcommand{\thetable}{S\arabic{table}}
\renewcommand{\thefigure}{S\arabic{figure}}
\renewcommand{\thesection}{S\arabic{section}}

\title{Supplementary information\\
  \large How much does preparation matter in molecular docking?}
\author{David Obeh Alobo}
\date{}

\begin{document}
\maketitle

All tables except Table~\ref{tab:si-protocols} are generated from the per-complex results
by \texttt{dockprepbench report}; the subset tables by the same command with
\texttt{--subset} \nolinkurl{configs/subsets/posebusters_308.txt}. Success is top-1 heavy-atom RMSD
below \SI{2}{\angstrom}; \pbvalid{} success also requires the top pose to pass every
PoseBusters redocking check other than its own RMSD check. Preparation failures count as failures. $\Delta$ is the change from the
same engine's baseline; gained and lost count complexes that succeed only in the arm or
only in the baseline; $p$ is the exact McNemar test, Holm-adjusted over the five arms of
each engine and metric.

\section{Success and paired comparisons for every engine}

\begin{table}[h]
  \centering
  \caption{Top-1 success for every engine and arm on all 428 complexes.}
  \label{tab:si-top1}
  \small
  \setlength{\tabcolsep}{3.5pt}
  \begin{tabular}{llrrrrr}
\toprule
Engine & Arm & Top-1 (\%) [95\% CI] & $\Delta$ (pts) [95\% CI] & Gained & Lost & Holm $p$ \\
\midrule
Vina & Baseline (4 \AA{} box) & 46.3 [41.6, 51.0] & -- & -- & -- & -- \\
 & Box padding 2 \AA{} & 53.0 [48.3, 57.7] & +6.8 [+3.5, +10.0] & 41 & 12 & $4.1\times10^{-4}$ \\
 & Box padding 6 \AA{} & 44.4 [39.8, 49.1] & $-1.9$ [$-4.3$, +0.5] & 10 & 18 & 0.37 \\
 & Box padding 8 \AA{} & 44.2 [39.5, 48.9] & $-2.1$ [$-5.2$, +1.0] & 19 & 28 & 0.37 \\
 & Protonation (pH 7.4) & 49.1 [44.4, 53.8] & +2.8 [$-0.1$, +5.7] & 26 & 14 & 0.24 \\
 & RDKit conformer & 50.7 [46.0, 55.4] & +4.4 [+0.2, +8.6] & 52 & 33 & 0.20 \\
\addlinespace
Smina & Baseline (4 \AA{} box) & 47.0 [42.3, 51.7] & -- & -- & -- & -- \\
 & Box padding 2 \AA{} & 54.7 [49.9, 59.3] & +7.7 [+4.4, +10.9] & 43 & 10 & $2.8\times10^{-5}$ \\
 & Box padding 6 \AA{} & 45.6 [40.9, 50.3] & $-1.4$ [$-3.8$, +1.0] & 11 & 17 & 0.69 \\
 & Box padding 8 \AA{} & 44.2 [39.5, 48.9] & $-2.8$ [$-5.8$, +0.2] & 16 & 28 & 0.38 \\
 & Protonation (pH 7.4) & 48.1 [43.4, 52.9] & +1.2 [$-1.5$, +3.8] & 19 & 14 & 0.69 \\
 & RDKit conformer & 50.2 [45.5, 54.9] & +3.3 [$-0.9$, +7.4] & 48 & 34 & 0.45 \\
\addlinespace
Gnina & Baseline (4 \AA{} box) & 62.6 [57.9, 67.1] & -- & -- & -- & -- \\
 & Box padding 2 \AA{} & 65.7 [61.0, 70.0] & +3.0 [$-0.5$, +6.5] & 36 & 23 & 0.47 \\
 & Box padding 6 \AA{} & 60.3 [55.6, 64.8] & $-2.3$ [$-5.7$, +1.0] & 22 & 32 & 0.47 \\
 & Box padding 8 \AA{} & 58.6 [53.9, 63.2] & $-4.0$ [$-7.8$, $-0.1$] & 28 & 45 & 0.30 \\
 & Protonation (pH 7.4) & 59.3 [54.6, 63.9] & $-3.3$ [$-7.2$, +0.7] & 30 & 44 & 0.47 \\
 & RDKit conformer & 63.3 [58.7, 67.7] & +0.7 [$-3.5$, +4.9] & 44 & 41 & 0.83 \\
\addlinespace
rDock & Baseline (4 \AA{} box) & 29.4 [25.3, 33.9] & -- & -- & -- & -- \\
 & Box padding 2 \AA{} & 41.1 [36.6, 45.8] & +11.7 [+7.9, +15.4] & 62 & 12 & $1.1\times10^{-8}$ \\
 & Box padding 6 \AA{} & 19.9 [16.4, 23.9] & $-9.6$ [$-12.9$, $-6.3$] & 7 & 48 & $3.9\times10^{-8}$ \\
 & Box padding 8 \AA{} & 17.3 [14.0, 21.2] & $-12.1$ [$-15.7$, $-8.6$] & 7 & 59 & $1.2\times10^{-10}$ \\
 & Protonation (pH 7.4) & 29.7 [25.5, 34.2] & +0.2 [$-3.1$, +3.6] & 27 & 26 & 1.0 \\
 & RDKit conformer & 25.5 [21.6, 29.8] & $-4.0$ [$-7.6$, $-0.4$] & 23 & 40 & 0.086 \\
\bottomrule
\end{tabular}

\end{table}

\begin{table}[h]
  \centering
  \caption{\pbvalid{} success for every engine and arm on all 428 complexes.}
  \label{tab:si-pbvalid}
  \small
  \setlength{\tabcolsep}{3.5pt}
  \begin{tabular}{llrrrrr}
\toprule
Engine & Arm & Pb-valid (\%) [95\% CI] & $\Delta$ (pts) [95\% CI] & Gained & Lost & Holm $p$ \\
\midrule
Vina & Baseline (4 \AA{} box) & 41.8 [37.2, 46.5] & -- & -- & -- & -- \\
 & Box padding 2 \AA{} & 47.9 [43.2, 52.6] & +6.1 [+2.7, +9.4] & 41 & 15 & 0.0034 \\
 & Box padding 6 \AA{} & 40.2 [35.6, 44.9] & $-1.6$ [$-4.2$, +0.9] & 12 & 19 & 0.56 \\
 & Box padding 8 \AA{} & 40.4 [35.9, 45.1] & $-1.4$ [$-4.6$, +1.8] & 21 & 27 & 0.56 \\
 & Protonation (pH 7.4) & 45.3 [40.7, 50.1] & +3.5 [+0.5, +6.5] & 30 & 15 & 0.14 \\
 & RDKit conformer & 45.3 [40.7, 50.1] & +3.5 [$-0.7$, +7.6] & 49 & 34 & 0.37 \\
\addlinespace
Smina & Baseline (4 \AA{} box) & 43.7 [39.1, 48.4] & -- & -- & -- & -- \\
 & Box padding 2 \AA{} & 49.3 [44.6, 54.0] & +5.6 [+2.5, +8.7] & 36 & 12 & 0.0036 \\
 & Box padding 6 \AA{} & 41.8 [37.2, 46.5] & $-1.9$ [$-4.4$, +0.6] & 11 & 19 & 0.60 \\
 & Box padding 8 \AA{} & 40.7 [36.1, 45.4] & $-3.0$ [$-5.9$, $-0.2$] & 13 & 26 & 0.21 \\
 & Protonation (pH 7.4) & 44.6 [40.0, 49.4] & +0.9 [$-1.9$, +3.7] & 21 & 17 & 0.63 \\
 & RDKit conformer & 46.3 [41.6, 51.0] & +2.6 [$-1.6$, +6.7] & 47 & 36 & 0.60 \\
\addlinespace
Gnina & Baseline (4 \AA{} box) & 55.8 [51.1, 60.5] & -- & -- & -- & -- \\
 & Box padding 2 \AA{} & 58.4 [53.7, 63.0] & +2.6 [$-1.2$, +6.3] & 40 & 29 & 0.77 \\
 & Box padding 6 \AA{} & 53.3 [48.5, 57.9] & $-2.6$ [$-6.1$, +0.9] & 24 & 35 & 0.77 \\
 & Box padding 8 \AA{} & 52.6 [47.8, 57.3] & $-3.3$ [$-7.2$, +0.7] & 31 & 45 & 0.68 \\
 & Protonation (pH 7.4) & 53.7 [49.0, 58.4] & $-2.1$ [$-5.9$, +1.7] & 31 & 40 & 0.77 \\
 & RDKit conformer & 55.4 [50.6, 60.0] & $-0.5$ [$-4.9$, +4.0] & 47 & 49 & 0.92 \\
\addlinespace
rDock & Baseline (4 \AA{} box) & 26.2 [22.2, 30.5] & -- & -- & -- & -- \\
 & Box padding 2 \AA{} & 35.5 [31.1, 40.2] & +9.3 [+5.7, +13.0] & 54 & 14 & $3.3\times10^{-6}$ \\
 & Box padding 6 \AA{} & 17.3 [14.0, 21.2] & $-8.9$ [$-12.2$, $-5.6$] & 8 & 46 & $5.5\times10^{-7}$ \\
 & Box padding 8 \AA{} & 15.7 [12.5, 19.4] & $-10.5$ [$-14.1$, $-7.0$] & 9 & 54 & $3.1\times10^{-8}$ \\
 & Protonation (pH 7.4) & 26.9 [22.9, 31.3] & +0.7 [$-2.8$, +4.2] & 30 & 27 & 0.79 \\
 & RDKit conformer & 22.4 [18.7, 26.6] & $-3.7$ [$-7.2$, $-0.3$] & 21 & 37 & 0.096 \\
\bottomrule
\end{tabular}

\end{table}

\section{Outcome changes against seed-to-seed variation}

\begin{table}[h]
  \centering
  \caption{Complexes whose top-1 outcome differs between each arm and its baseline. For
    Vina, whose baseline was repeated with four seeds, the ratio is to the mean number of
    complexes that differ between two seeds (17.5).}
  \label{tab:si-flips}
  \small
  \begin{tabular}{llrrr}
\toprule
Engine & Arm & $\Delta$ (pts) & Flipped & Ratio to seed flips \\
\midrule
Vina & Box padding 2 \AA{} & +6.8 & 53 & 3.0 \\
 & Box padding 6 \AA{} & $-1.9$ & 28 & 1.6 \\
 & Box padding 8 \AA{} & $-2.1$ & 47 & 2.7 \\
 & Protonation (pH 7.4) & +2.8 & 40 & 2.3 \\
 & RDKit conformer & +4.4 & 85 & 4.9 \\
\addlinespace
Smina & Box padding 2 \AA{} & +7.7 & 53 & -- \\
 & Box padding 6 \AA{} & $-1.4$ & 28 & -- \\
 & Box padding 8 \AA{} & $-2.8$ & 44 & -- \\
 & Protonation (pH 7.4) & +1.2 & 33 & -- \\
 & RDKit conformer & +3.3 & 82 & -- \\
\addlinespace
Gnina & Box padding 2 \AA{} & +3.0 & 59 & -- \\
 & Box padding 6 \AA{} & $-2.3$ & 54 & -- \\
 & Box padding 8 \AA{} & $-4.0$ & 73 & -- \\
 & Protonation (pH 7.4) & $-3.3$ & 74 & -- \\
 & RDKit conformer & +0.7 & 85 & -- \\
\bottomrule
\end{tabular}

\end{table}

\clearpage
\section{The 308-complex subset}

The journal version of the PoseBusters study reports results on 308 of the 428 complexes,
leaving out those whose ligand lies within \SI{5}{\angstrom} of a symmetry-related copy of
the protein (a crystal contact; Methods). Tables~\ref{tab:si-308-top1}
and~\ref{tab:si-308-pbvalid} repeat Tables~\ref{tab:si-top1} and~\ref{tab:si-pbvalid} on
that subset; Holm adjustment is recomputed within the subset.

\begin{table}[h]
  \centering
  \caption{Top-1 success for every engine and arm on the 308-complex subset.}
  \label{tab:si-308-top1}
  \small
  \setlength{\tabcolsep}{3.5pt}
  \begin{tabular}{llrrrrr}
\toprule
Engine & Arm & Top-1 (\%) [95\% CI] & $\Delta$ (pts) & Gained & Lost & Holm $p$ \\
\midrule
Vina & Baseline (4 \AA{} box) & 52.9 [47.3, 58.4] & -- & -- & -- & -- \\
 & Box padding 2 \AA{} & 58.4 [52.9, 63.8] & +5.5 & 25 & 8 & 0.023 \\
 & Box padding 6 \AA{} & 50.3 [44.8, 55.9] & $-2.6$ & 4 & 12 & 0.31 \\
 & Box padding 8 \AA{} & 49.7 [44.1, 55.2] & $-3.2$ & 9 & 19 & 0.31 \\
 & Protonation (pH 7.4) & 54.5 [49.0, 60.0] & +1.6 & 17 & 12 & 0.46 \\
 & RDKit conformer & 56.5 [50.9, 61.9] & +3.6 & 34 & 23 & 0.37 \\
\addlinespace
Smina & Baseline (4 \AA{} box) & 53.9 [48.3, 59.4] & -- & -- & -- & -- \\
 & Box padding 2 \AA{} & 60.1 [54.5, 65.4] & +6.2 & 29 & 10 & 0.017 \\
 & Box padding 6 \AA{} & 52.3 [46.7, 57.8] & $-1.6$ & 7 & 12 & 1.0 \\
 & Box padding 8 \AA{} & 50.3 [44.8, 55.9] & $-3.6$ & 8 & 19 & 0.21 \\
 & Protonation (pH 7.4) & 53.9 [48.3, 59.4] & +0.0 & 11 & 11 & 1.0 \\
 & RDKit conformer & 56.5 [50.9, 61.9] & +2.6 & 33 & 25 & 1.0 \\
\addlinespace
Gnina & Baseline (4 \AA{} box) & 68.8 [63.5, 73.7] & -- & -- & -- & -- \\
 & Box padding 2 \AA{} & 71.8 [66.5, 76.5] & +2.9 & 25 & 16 & 0.63 \\
 & Box padding 6 \AA{} & 66.2 [60.8, 71.3] & $-2.6$ & 16 & 24 & 0.63 \\
 & Box padding 8 \AA{} & 63.6 [58.1, 68.8] & $-5.2$ & 18 & 34 & 0.15 \\
 & Protonation (pH 7.4) & 62.7 [57.1, 67.9] & $-6.2$ & 19 & 38 & 0.082 \\
 & RDKit conformer & 69.8 [64.5, 74.7] & +1.0 & 32 & 29 & 0.80 \\
\bottomrule
\end{tabular}

\end{table}

\begin{table}[h]
  \centering
  \caption{\pbvalid{} success for every engine and arm on the 308-complex subset.}
  \label{tab:si-308-pbvalid}
  \small
  \setlength{\tabcolsep}{3.5pt}
  \begin{tabular}{llrrrrr}
\toprule
Engine & Arm & Pb-valid (\%) [95\% CI] & $\Delta$ (pts) & Gained & Lost & Holm $p$ \\
\midrule
Vina & Baseline (4 \AA{} box) & 47.4 [41.9, 53.0] & -- & -- & -- & -- \\
 & Box padding 2 \AA{} & 52.3 [46.7, 57.8] & +4.9 & 25 & 10 & 0.083 \\
 & Box padding 6 \AA{} & 44.8 [39.3, 50.4] & $-2.6$ & 5 & 13 & 0.39 \\
 & Box padding 8 \AA{} & 45.1 [39.7, 50.7] & $-2.3$ & 11 & 18 & 0.49 \\
 & Protonation (pH 7.4) & 50.3 [44.8, 55.9] & +2.9 & 21 & 12 & 0.49 \\
 & RDKit conformer & 50.6 [45.1, 56.2] & +3.2 & 33 & 23 & 0.49 \\
\addlinespace
Smina & Baseline (4 \AA{} box) & 50.0 [44.5, 55.5] & -- & -- & -- & -- \\
 & Box padding 2 \AA{} & 52.9 [47.3, 58.4] & +2.9 & 21 & 12 & 0.65 \\
 & Box padding 6 \AA{} & 48.1 [42.5, 53.6] & $-1.9$ & 8 & 14 & 0.86 \\
 & Box padding 8 \AA{} & 45.8 [40.3, 51.4] & $-4.2$ & 5 & 18 & 0.053 \\
 & Protonation (pH 7.4) & 49.7 [44.1, 55.2] & $-0.3$ & 13 & 14 & 1.0 \\
 & RDKit conformer & 52.3 [46.7, 57.8] & +2.3 & 34 & 27 & 0.89 \\
\addlinespace
Gnina & Baseline (4 \AA{} box) & 61.4 [55.8, 66.6] & -- & -- & -- & -- \\
 & Box padding 2 \AA{} & 63.3 [57.8, 68.5] & +1.9 & 29 & 23 & 0.98 \\
 & Box padding 6 \AA{} & 58.4 [52.9, 63.8] & $-2.9$ & 18 & 27 & 0.70 \\
 & Box padding 8 \AA{} & 56.8 [51.2, 62.2] & $-4.5$ & 20 & 34 & 0.31 \\
 & Protonation (pH 7.4) & 56.5 [50.9, 61.9] & $-4.9$ & 21 & 36 & 0.31 \\
 & RDKit conformer & 60.4 [54.8, 65.7] & $-1.0$ & 33 & 36 & 0.98 \\
\bottomrule
\end{tabular}

\end{table}

\clearpage
\section{Comparison with the PoseBusters Vina protocol}

\begin{table}[h]
  \centering
  \caption{Vina protocols of the PoseBusters study \citep{buttenschoen2024posebusters} and
    of this work's baseline, on the same 428 complexes. Per-complex results of the
    PoseBusters run are from its Zenodo record \citep{posebusters2023data}; the comparison
    is reproduced by \texttt{scripts/compare\_posebusters\_vina.py}. PoseBusters settings
    are from section S1 of that study's electronic supplementary information.}
  \label{tab:si-protocols}
  \small
  \begin{tabular}{lll}
    \toprule
    & PoseBusters & This work (baseline) \\
    \midrule
    Vina version & 1.2.3 & 1.2.7 \\
    Exhaustiveness & 32 & 8 \\
    Poses & 40 & 9 \\
    Search box & \SI{25}{\angstrom} cube on ligand centre & ligand extent + \SI{4}{\angstrom} per side \\
    Receptor & Reduce, ADFR & PDBFixer, OpenMM, Meeko \\
    Ligand start & RDKit ETKDGv3 + UFF & supplied (same set) \\
    \midrule
    Top-1 success & 52.3\% (224/428) & 46.3\% (198/428) \\
    \pbvalid{} success & 51.2\% (219/428) & 41.8\% (179/428) \\
    \bottomrule
  \end{tabular}
\end{table}

\section{The replaced rDock runs}

rDock defines its search space as a cavity rather than a box, and two cavity definitions
were tried and replaced before the final one (Table~\ref{tab:si-rdock-versions}); their
results are kept for the record.

The first cavity was a sphere centred on the search box and passing through its corners,
chosen so that it enclosed the whole box. That sphere is much larger than the box, and the
run was compromised: baseline top-1 success was 13.6\% (58/428), 129 top poses lay in the
wrong site (Vina: 20), and 100 complexes failed, many because their docking processes were
stopped by the memory limit of a 64-worker batch while others returned poses RDKit could
not read (rDock writes poses without formal charges).

The second cavity was mapped around the crystal ligand, keeping free space within the box
padding of any ligand atom, and each pose's chemistry was restored from the input ligand.
Its success was implausibly high, 42.5\% at baseline and 72.9\% with
\SI{2}{\angstrom} padding, well above Gnina: with a small padding the cavity takes the shape
of the crystal ligand itself, so it encodes the answer rather than the search space the
other engines were given.

The final cavity (Methods) is the free space inside the same box the other engines use,
mapped from a grid of points filling the box, with the same chemistry restoration. It keeps
the search space comparable across engines without revealing the ligand's shape.

\begin{table}[h]
  \centering
  \caption{rDock top-1 success on all 428 complexes under the three cavity definitions.
    Failures count as failures.}
  \label{tab:si-rdock-versions}
  \small
  \begin{tabular}{lrrr}
    \toprule
    Cavity & Baseline (\%) & Padding \SI{2}{\angstrom} (\%) & Failed at baseline \\
    \midrule
    Sphere enclosing the box (replaced) & 13.6 & 18.5 & 100 \\
    Around the crystal ligand (replaced) & 42.5 & 72.9 & 0 \\
    Free space in the box (final) & 29.4 & 41.1 & 0 \\
    \bottomrule
  \end{tabular}
\end{table}

\clearpage
\section{Which complexes each arm affects}

For each engine and arm, complexes whose top-1 outcome differs from the baseline
(flipped) are compared with the rest. For continuous properties the table gives medians
and the rank-biserial correlation $r$ (Mann--Whitney $U$ test); for binary properties the
share of complexes with the property and the odds ratio (Fisher's exact test). P-values
are Holm-adjusted over the properties of each engine and arm. Seed sensitivity is
available for Vina only.

\newcommand{\SensitivityCaption}{Ligand and complex properties of flipped and unflipped
  complexes, all 428 complexes.}
{\footnotesize\begin{longtable}{lllrrrr}
\caption{\SensitivityCaption}\label{tab:si-sensitivity}\\
\toprule
Engine & Arm & Property & Flipped & Unflipped & Effect & Holm $p$ \\
\midrule
\endfirsthead
\toprule
Engine & Arm & Property & Flipped & Unflipped & Effect & Holm $p$ \\
\midrule
\endhead
\bottomrule
\endfoot
Vina & Box padding 2 \AA{} & Heavy atoms & 29 & 24 & $r = 0.21$ & 0.057 \\
 &  & Rotatable bonds & 7 & 6 & $r = 0.18$ & 0.15 \\
 &  & Absolute charge & 0 & 0 & $r = -0.01$ & 1.0 \\
 &  & Charge changed at pH 7.4 & 62.3\% & 63.5\% & OR 0.95 & 1.0 \\
 &  & Cofactor present & 50.9\% & 57.9\% & OR 0.76 & 1.0 \\
 &  & Seed-sensitive & 35.8\% & 3.7\% & OR 14 & $3.3\times10^{-10}$ \\
\addlinespace
Vina & Box padding 6 \AA{} & Heavy atoms & 26.5 & 25 & $r = 0.16$ & 0.71 \\
 &  & Rotatable bonds & 7 & 6 & $r = 0.17$ & 0.71 \\
 &  & Absolute charge & 0 & 0 & $r = -0.03$ & 1.0 \\
 &  & Charge changed at pH 7.4 & 71.4\% & 62.7\% & OR 1.5 & 1.0 \\
 &  & Cofactor present & 53.6\% & 57.2\% & OR 0.86 & 1.0 \\
 &  & Seed-sensitive & 60.7\% & 4.0\% & OR 37 & $9.0\times10^{-14}$ \\
\addlinespace
Vina & Box padding 8 \AA{} & Heavy atoms & 24 & 25 & $r = -0.07$ & 1.0 \\
 &  & Rotatable bonds & 6 & 6 & $r = -0.07$ & 1.0 \\
 &  & Absolute charge & 0 & 0 & $r = -0.03$ & 1.0 \\
 &  & Charge changed at pH 7.4 & 61.7\% & 63.5\% & OR 0.93 & 1.0 \\
 &  & Cofactor present & 51.1\% & 57.7\% & OR 0.76 & 1.0 \\
 &  & Seed-sensitive & 40.4\% & 3.7\% & OR 18 & $2.4\times10^{-11}$ \\
\addlinespace
Vina & Protonation (pH 7.4) & Heavy atoms & 29 & 25 & $r = 0.18$ & 0.18 \\
 &  & Rotatable bonds & 8 & 6 & $r = 0.19$ & 0.18 \\
 &  & Absolute charge & 0 & 0 & $r = 0.03$ & 0.49 \\
 &  & Charge changed at pH 7.4 & 80.0\% & 61.6\% & OR 2.5 & 0.12 \\
 &  & Cofactor present & 55.0\% & 57.2\% & OR 0.91 & 0.87 \\
 &  & Seed-sensitive & 40.0\% & 4.4\% & OR 15 & $3.6\times10^{-9}$ \\
\addlinespace
Vina & RDKit conformer & Heavy atoms & 25 & 25 & $r = 0.01$ & 1.0 \\
 &  & Rotatable bonds & 6 & 6 & $r = 0.00$ & 1.0 \\
 &  & Absolute charge & 0 & 0 & $r = 0.00$ & 1.0 \\
 &  & Charge changed at pH 7.4 & 58.8\% & 64.4\% & OR 0.79 & 1.0 \\
 &  & Cofactor present & 54.1\% & 57.7\% & OR 0.86 & 1.0 \\
 &  & Seed-sensitive & 25.9\% & 3.2\% & OR 11 & $6.1\times10^{-9}$ \\
\addlinespace
Smina & Box padding 2 \AA{} & Heavy atoms & 28 & 24 & $r = 0.18$ & 0.13 \\
 &  & Rotatable bonds & 10 & 6 & $r = 0.31$ & 0.0012 \\
 &  & Absolute charge & 0 & 0 & $r = 0.02$ & 1.0 \\
 &  & Charge changed at pH 7.4 & 66.0\% & 62.9\% & OR 1.1 & 1.0 \\
 &  & Cofactor present & 56.6\% & 57.1\% & OR 0.98 & 1.0 \\
\addlinespace
Smina & Box padding 6 \AA{} & Heavy atoms & 27.5 & 25 & $r = 0.01$ & 1.0 \\
 &  & Rotatable bonds & 7 & 6 & $r = 0.07$ & 1.0 \\
 &  & Absolute charge & 0 & 0 & $r = 0.01$ & 1.0 \\
 &  & Charge changed at pH 7.4 & 75.0\% & 62.5\% & OR 1.8 & 1.0 \\
 &  & Cofactor present & 67.9\% & 56.2\% & OR 1.6 & 1.0 \\
\addlinespace
Smina & Box padding 8 \AA{} & Heavy atoms & 25.5 & 25 & $r = 0.03$ & 1.0 \\
 &  & Rotatable bonds & 6 & 6 & $r = 0.01$ & 1.0 \\
 &  & Absolute charge & 0 & 0 & $r = -0.03$ & 1.0 \\
 &  & Charge changed at pH 7.4 & 65.9\% & 63.0\% & OR 1.1 & 1.0 \\
 &  & Cofactor present & 61.4\% & 56.5\% & OR 1.2 & 1.0 \\
\addlinespace
Smina & Protonation (pH 7.4) & Heavy atoms & 23 & 25 & $r = -0.10$ & 1.0 \\
 &  & Rotatable bonds & 6 & 6 & $r = 0.09$ & 1.0 \\
 &  & Absolute charge & 0 & 0 & $r = 0.01$ & 1.0 \\
 &  & Charge changed at pH 7.4 & 72.7\% & 62.5\% & OR 1.6 & 1.0 \\
 &  & Cofactor present & 66.7\% & 56.2\% & OR 1.6 & 1.0 \\
\addlinespace
Smina & RDKit conformer & Heavy atoms & 26 & 25 & $r = 0.05$ & 1.0 \\
 &  & Rotatable bonds & 7.5 & 6 & $r = 0.13$ & 0.30 \\
 &  & Absolute charge & 0 & 0 & $r = -0.01$ & 1.0 \\
 &  & Charge changed at pH 7.4 & 62.2\% & 63.6\% & OR 0.94 & 1.0 \\
 &  & Cofactor present & 54.9\% & 57.5\% & OR 0.9 & 1.0 \\
\addlinespace
Gnina & Box padding 2 \AA{} & Heavy atoms & 26 & 25 & $r = 0.07$ & 1.0 \\
 &  & Rotatable bonds & 8 & 6 & $r = 0.17$ & 0.18 \\
 &  & Absolute charge & 0 & 0 & $r = 0.01$ & 1.0 \\
 &  & Charge changed at pH 7.4 & 67.8\% & 62.6\% & OR 1.3 & 1.0 \\
 &  & Cofactor present & 62.7\% & 56.1\% & OR 1.3 & 1.0 \\
\addlinespace
Gnina & Box padding 6 \AA{} & Heavy atoms & 22 & 25 & $r = -0.06$ & 1.0 \\
 &  & Rotatable bonds & 7 & 6 & $r = 0.05$ & 1.0 \\
 &  & Absolute charge & 0 & 0 & $r = 0.04$ & 0.49 \\
 &  & Charge changed at pH 7.4 & 68.5\% & 62.6\% & OR 1.3 & 1.0 \\
 &  & Cofactor present & 55.6\% & 57.2\% & OR 0.93 & 1.0 \\
\addlinespace
Gnina & Box padding 8 \AA{} & Heavy atoms & 24 & 25 & $r = 0.00$ & 0.98 \\
 &  & Rotatable bonds & 6 & 6 & $r = 0.09$ & 0.97 \\
 &  & Absolute charge & 0 & 0 & $r = 0.02$ & 0.97 \\
 &  & Charge changed at pH 7.4 & 71.2\% & 61.7\% & OR 1.5 & 0.71 \\
 &  & Cofactor present & 63.0\% & 55.8\% & OR 1.4 & 0.97 \\
\addlinespace
Gnina & Protonation (pH 7.4) & Heavy atoms & 24 & 25 & $r = -0.02$ & 1.0 \\
 &  & Rotatable bonds & 8 & 6 & $r = 0.16$ & 0.13 \\
 &  & Absolute charge & 0 & 0 & $r = 0.02$ & 0.87 \\
 &  & Charge changed at pH 7.4 & 73.0\% & 61.3\% & OR 1.7 & 0.26 \\
 &  & Cofactor present & 59.5\% & 56.5\% & OR 1.1 & 1.0 \\
\addlinespace
Gnina & RDKit conformer & Heavy atoms & 24 & 25 & $r = -0.02$ & 1.0 \\
 &  & Rotatable bonds & 7 & 6 & $r = 0.09$ & 1.0 \\
 &  & Absolute charge & 0 & 0 & $r = 0.01$ & 1.0 \\
 &  & Charge changed at pH 7.4 & 69.4\% & 61.8\% & OR 1.4 & 1.0 \\
 &  & Cofactor present & 56.5\% & 57.1\% & OR 0.97 & 1.0 \\
\addlinespace
\end{longtable}
}

\bibliographystyle{unsrtnat}
\bibliography{references}

\end{document}
